\subsection{紧空间的积}

定理 3（Tychonoff）. 拟紧（相应地，紧）空间的每个积都是拟紧（相应地，紧）的。反之，若非空空间的积是拟紧（相应地，紧）的，则每个因子都是拟紧（相应地，紧）的。

鉴于 § 8 第 2 目命题 7 给出的豪斯多夫积空间的刻画，只需对拟紧空间证明这些断言。若 \(\mathbf{X} = \prod_{\iota \in I} \mathbf{X}_{\iota}\) 是拟紧的且非空，则由第 4 目定理 2，\(\mathbf{X}_{\iota}=\mathrm{pr}_{\iota}(\mathbf{X})\) 是拟紧的。反之，设诸 \(X_{\iota}\) 是拟紧的，并设 U 是 X 上的一个超滤子；则对每个 \(\iota\in I\)，\(\mathrm{pr}_{\iota}(\mathfrak{U})\) 是 \(X_{\iota}\) 上的一个超滤子基（§ 6，第 6 目，命题 10），由公理 (C') 它收敛；于是 U 收敛（§ 7，第 6 目，命题 10 的推论 1），因此 X 是拟紧的。

推论. 拓扑空间之积的一个子集是相对拟紧的，必要且充分的是，它的每个投影在相应的因子中都是相对拟紧的。

必要性由 § 4 定理 2 得出。为证充分性，设 A 是 \(\prod_{i} X_{i}\) 的一个子集，使得对每个指标 \(i\)，\(\operatorname{pr}_{i}(A)\) 包含于拟紧子集 \(K_{i}\)（\(X_{i}\) 的）中；则 A 包含于拟紧子集 \(\prod_{i} K_{i}\)（\(\prod_{i} X_{i}\) 的）中。

